\chapter{Introduction}

Large Language Models (LLMs) have transformed the field of natural
language processing, enabling computational systems to generate coherent, contextually
relevant responses across a wide range of tasks.
However, LLMs are constrained by a fundamental limitation: their knowledge is fixed
at the point of training.
When queried on topics beyond this boundary, models frequently produce responses that
are factually incorrect, a behaviour commonly referred to as
hallucination.

Retrieval-Augmented Generation (RAG) was introduced to mitigate this limitation by
enabling models to retrieve relevant information from external sources before generating
a response.
This approach grounds model outputs in real-world data and reduces reliance on static
parametric knowledge.
Subsequent work extended this paradigm further — first by demonstrating that fusing
multiple retrieved passages yields measurable gains on open-domain question
answering, and then by integrating retrieval directly into the pre-training stage,
producing representations that are inherently more
knowledge-aware.

Despite these advances, standard RAG implementations operate through a linear,
single-pass pipeline that retrieves a document, reads it once, and produces a single
response.
This design assumes a single valid interpretation of the retrieved content, which is
insufficient for complex analytical tasks that require critical evaluation,
cross-verification, or multi-perspective reasoning.
Efforts to address this limitation have led to frameworks that interleave retrieval
with iterative reasoning, and to self-reflective generation
mechanisms that allow a model to evaluate the quality of both its retrieved evidence
and its generated output before producing a final answer.

In addition, most high-performance RAG systems depend on cloud-based infrastructure,
requiring user queries and potentially sensitive data to be transmitted to external
servers.
This raises significant concerns regarding data privacy and sovereignty, particularly
in professional, medical, or research contexts where confidentiality is
essential.
Furthermore, existing systems are largely static; they do not adapt based on past
interactions or evolving user requirements, making repeated errors likely and limiting
their long-term utility.

To address these limitations, this project proposes \textbf{EvoRAG (Evolutionary
Multi-Persona RAG)}, a locally deployed, multi-agent framework that enhances
retrieval-augmented generation through structured cognitive
diversity.
EvoRAG orchestrates eight distinct AI personas, each with a different reasoning style,
to analyse retrieved content simultaneously and independently.
Their outputs are evaluated through an internal consensus mechanism that scores
responses based on factual accuracy, analytical depth, and clarity.
The system further incorporates an \emph{Evolutionary Learning Engine} that monitors
persona performance over time and autonomously replaces underperforming agents with
improved variants, based on internal win rates and user feedback.
All processing is performed on local hardware, ensuring that no user data is transmitted
externally.

% ------------------------------------------------------------------
\section{Motivation}

The decision to undertake EvoRAG as a project was driven by a confluence of
real-world observations and unresolved technical gaps in the current landscape of
AI-driven information systems.
As large language models became embedded in decision-making across healthcare, education,
legal services, and public discourse, a clear pattern emerged: these systems are capable
of generating highly confident yet factually incorrect responses, a behaviour with direct
consequences for anyone relying on them.
Witnessing the gap between the trust users place in AI-generated answers and the actual
reliability of those answers made it evident that the existing single-pass retrieval and
generation paradigm was insufficient for high-stakes analytical tasks.

A further motivating observation was the structural incompatibility of cloud-dependent
RAG deployments with privacy-sensitive domains.
In fields such as clinical medicine, legal practice, and enterprise research, routing
queries through external servers is not merely inconvenient — it may constitute a
direct regulatory violation.
The absence of a capable, locally deployable alternative that could match the analytical
depth of hosted systems represented a concrete and pressing gap that this project
is designed to fill.

The project was also motivated by the absence of any mechanism for transparency and
explainability in how existing RAG systems arrive at their responses.
A single answer, produced by a single reasoning trajectory, provides no insight into
what alternative interpretations were considered or rejected.
For users in professional contexts, this opacity undermines trust and limits
accountability.
EvoRAG was conceived to address this by making the diversity of reasoning visible —
eight independent personas each produce an analytical response before a consensus
is formed, giving users a richer and more auditable view of the reasoning process.

On the technical side, the specific limitations of existing RAG architectures provided
equally compelling motivation.
Standard pipelines retrieve a document, process it once, and emit a single output with
no internal mechanism to challenge, verify, or cross-validate the result.
Recent work on retrieval-interleaved reasoning and self-reflective generation has
shown the value of iterative evaluation, yet
these advances have not been integrated into a unified, locally deployable,
consensus-driven framework.
Furthermore, all deployed RAG systems observed during the course of this project shared
a common architectural flaw: they do not improve over time.
Query history, user feedback, and observed performance outcomes are discarded after
each session, making systematic error correction impossible without manual retraining.
This motivated the design of EvoRAG's evolutionary learning engine, which autonomously
refines persona configurations based on empirical win rates and accumulated feedback,
enabling the system to become more effective the longer it is used.

% ------------------------------------------------------------------
\section{Objectives of the Project}
% ------------------------------------------------------------------

The primary goals of the EvoRAG project are:

\begin{itemize}
    \item To build a robust pipeline that can autonomously retrieve and index external
          knowledge for local use at runtime, in direct response to each user query.
    \item To design a multi-agent framework that simulates eight distinct AI personas
          to provide diverse, parallel analytical views of the same retrieved context.
    \item To implement a peer-review scoring system that programmatically filters out
          weak or inaccurate responses through a structured consensus mechanism.
    \item To develop an evolutionary tracking system that improves the system's
          reasoning capability over time based on win rates and user feedback, without
          manual reconfiguration or retraining of the base model.
    \item To ensure all operations are local and privacy-preserving, removing the
          dependency on external APIs or cloud-based services entirely.
\end{itemize}

\section{Organisation of the Report}

This report is structured to provide comprehensive coverage of the proposed system's
design and implementation. The content is organised across five chapters as follows:

\begin{itemize}
    \item \textbf{Chapter 1: Introduction} — Presents the background, motivation, and
          objectives of the EvoRAG project, establishing the research context and the
          limitations of existing RAG systems that this work seeks to address.

    \item \textbf{Chapter 2: Literature Survey} — Surveys the relevant literature
          across agentic RAG frameworks, multi-agent collaboration, memory management,
          and evaluation methodology, and identifies the specific research gaps that
          EvoRAG is designed to address.

    \item \textbf{Chapter 3: System Overview} — Describes the functional and
          non-functional requirements of the system alongside the hardware, software,
          and overall architectural design of the proposed solution.

    \item \textbf{Chapter 4: Methodology} — Details the implementation approach,
          covering dataset preparation, knowledge ingestion, and the algorithmic
          operation of each core engine within the EvoRAG pipeline.

    \item \textbf{Chapter 5: Conclusion} — Presents the overall conclusions of the
          project, reflects on current limitations of the implemented system, and
          outlines directions for future research and development.
\end{itemize}